% Adaptación de la plantilla create-pdf para la memoria de la práctica 2.
\usepackage[a4paper,left=2.25cm,right=2.25cm,top=2.7cm,bottom=2.5cm]{geometry}
\usepackage{setspace}
\setstretch{1.13}
\usepackage{graphicx,booktabs,array,tabularx,float,caption}
\usepackage{fancyhdr,titlesec,lastpage,needspace,etoolbox,xurl}
\usepackage{fvextra}
\usepackage{newunicodechar}
\newunicodechar{−}{\ensuremath{-}}
\newunicodechar{≤}{\ensuremath{\leq}}
\newunicodechar{≥}{\ensuremath{\geq}}
\newunicodechar{→}{\ensuremath{\rightarrow}}
\newunicodechar{×}{\ensuremath{\times}}
\newunicodechar{⁹}{\ensuremath{^9}}
\definecolor{uieblue}{HTML}{315F9B}
\definecolor{inkgray}{HTML}{555E6A}
\titleformat{\section}{\Large\bfseries\color{uieblue}}{\thesection.}{0.55em}{}
\titleformat{\subsection}{\large\bfseries}{\thesubsection.}{0.55em}{}
\titlespacing*{\section}{0pt}{20pt}{9pt}
\titlespacing*{\subsection}{0pt}{14pt}{6pt}
\pretocmd{\section}{\Needspace{8\baselineskip}}{}{}
\setlength{\parindent}{0pt}
\setlength{\parskip}{5pt}
\setlength{\headheight}{30pt}
\setlength{\headsep}{14pt}
\pagestyle{fancy}
\fancyhf{}
\fancyhead[L]{\includegraphics[height=0.8cm]{assets/logo_uie_mark.png}}
\fancyhead[R]{\small\textcolor{inkgray}{Gemelos Digitales\enspace |\enspace Práctica 2}}
\fancyfoot[L]{\footnotesize\textcolor{inkgray}{Alan Salazar · UIE}}
\fancyfoot[C]{\footnotesize\textcolor{inkgray}{30 de septiembre de 2026}}
\fancyfoot[R]{\footnotesize\thepage\ /\ \pageref*{LastPage}}
\renewcommand{\headrulewidth}{0.4pt}
\renewcommand{\footrulewidth}{0.3pt}
\fancypagestyle{plain}{\pagestyle{fancy}}
\captionsetup{font=small,labelfont=bf,labelsep=period,justification=raggedright,singlelinecheck=false,skip=7pt}
\addto\captionsspanish{\renewcommand{\figurename}{Figura}\renewcommand{\tablename}{Tabla}\renewcommand{\contentsname}{Índice}}
\AtBeginEnvironment{longtable}{\small\setstretch{1.04}}
\setlength{\tabcolsep}{4pt}
\renewcommand{\arraystretch}{1.15}
\fvset{breaklines=true,breakanywhere=true,fontsize=\footnotesize}
\RecustomVerbatimEnvironment{verbatim}{Verbatim}{breaklines=true,breakanywhere=true,fontsize=\footnotesize}
\setcounter{tocdepth}{1}
\setlength{\emergencystretch}{3em}
\clubpenalty=10000
\widowpenalty=10000
\AtBeginDocument{
\hypersetup{pdftitle={Práctica 2: Flota de Helsinki por MQTT},pdfauthor={Alan Salazar},pdfsubject={Memoria técnica de Gemelos Digitales · UIE}}
\thispagestyle{fancy}
\vspace*{0.8cm}
\begin{center}
\includegraphics[width=0.88\linewidth]{assets/logo_uie.png}\\[1.2cm]
{\large Grado en Ingeniería en Sistemas Inteligentes}\\[0.25cm]
{\normalsize Cuarto curso · Primer semestre}\\[0.6cm]
{\Large\color{uieblue} Gemelos Digitales}\\[1.0cm]
{\Huge\bfseries\color{uieblue} Flota de Helsinki\\[0.15cm] por MQTT}\\[0.45cm]
{\Large Práctica 2 · Memoria técnica}\\[0.8cm]
\rule{0.85\linewidth}{0.5pt}\\[0.45cm]
\begin{minipage}{0.86\linewidth}
\textbf{Objetivo de la entrega}\\[0.2cm]
Explicar la implementación de un sistema que recibe datos de tranvías, valida y almacena sus medidas y muestra su estado. La memoria desarrolla también la decisión del plan de invierno y la evaluación de la predicción de llegada.
\end{minipage}\\[0.45cm]
\rule{0.85\linewidth}{0.5pt}\\[0.7cm]
{\large Alan Salazar}\\[0.25cm]
{\normalsize Universidad Intercontinental de la Empresa}\\[0.25cm]
{\normalsize 30 de septiembre de 2026}
\end{center}
\clearpage
\tableofcontents
\clearpage
}
